% ECONOMICS_PROBLEM: {"id": "E52-261", "title": "Optimal inflation target", "jel_code": "E52", "source_id": "OP-0247"}
% !TeX program = xelatex
\documentclass[11pt,a4paper]{article}
\usepackage[margin=23mm]{geometry}
\usepackage{fontspec}
\setmainfont{Latin Modern Roman}
\usepackage{amsmath,amssymb,mathtools}
\usepackage{xcolor,enumitem,booktabs,longtable,array,xurl,hyperref}
\definecolor{muted}{HTML}{53636C}
\hypersetup{colorlinks=true,urlcolor=blue,linkcolor=blue}
\setlength{\parindent}{0pt}
\setlength{\parskip}{5pt}
\newcommand{\R}{\mathbb R}
\newcommand{\E}{\mathbb E}
\newcommand{\N}{\mathbb N}
\newcommand{\ind}{\mathbf 1}
\newcommand{\Var}{\operatorname{Var}}
\newcommand{\Cov}{\operatorname{Cov}}
\newcommand{\supp}{\operatorname{supp}}
\DeclareMathOperator*{\argmin}{arg\,min}
\DeclareMathOperator*{\argmax}{arg\,max}
\newcommand{\entrymeta}[3]{{\sffamily\small\color{muted}\textbf{\texttt{#1}}\quad|\quad #2\par #3\par}}
\newcommand{\Problemref}[1]{\texttt{#1}}
\newcommand{\JELref}[2]{\texttt{#2} (JEL \texttt{#1})}
\begin{document}
\section*{Optimal inflation target}
\textbf{Problem ID:} \texttt{E52-261}\quad \textbf{JEL:} \texttt{E52}\par
% BEGIN ECONOMICS STATEMENT
\label{problem:OP-0247}
{\sffamily\small\color{muted}Original subject group: Inflation and monetary policy\par}
\label{definitions:problem:30007528}
\textbf{Audit classification:} Background; proposed extension. Strategy comparisons do not publish an unresolved optimal numerical target; the robust-target optimization is editorial. Verification concerns the cited version, not first priority or proof that this exact editorial specification remains unsolved on October 8, 2026.
\subsection*{Definitions and precise research target}
\paragraph{Editorial mathematical specification.} Fix a fully specified stochastic monetary economy with household types \(i\), social weights \(\omega_i\), sticky prices, downward nominal wage constraints, and an effective lower bound \(i_t\ge\underline i\) on nominal policy rates. Preferences \(u_i(c,\ell)\), shock laws, fiscal financing, and the admissible class \(\mathcal R(\bar\pi)\) of policy rules around a constant inflation target \(\bar\pi\) are given. Restrict targets to a compact interval \([\pi_L,\pi_U]\) and require a declared equilibrium selection and integrable discounted utility for every evaluated rule. Here \(\beta\in(0,1)\), and nonnegative welfare weights sum to one; assume a maximizing target exists or report only the supremum.
Define
\[
W(\bar\pi)=\sup_{r\in\mathcal R(\bar\pi)}
E_r\sum_{t\ge0}\beta^t\sum_i\omega_i u_i(c_{it},\ell_{it}),\qquad
\bar\pi^*\in\arg\max_{\bar\pi\in[\pi_L,\pi_U]}W(\bar\pi).
\]
Determine an empirically credible set for \(\bar\pi^*\) when wage rigidity, neutral-rate dynamics, and household inflation exposure are only partially identified. Quantify the tradeoff between lower-bound stabilization benefits, price dispersion, and heterogeneous purchasing-power costs within the same welfare criterion. Report whether the ranking of targets is robust to admissible parameter uncertainty and policy-rule restrictions. An optimum calculated in one calibrated model is a conditional answer, not a universal resolution. This specification asks for a defensible robust target in a declared environment; social weights and institutional policy choices cannot be inferred solely from inflation data.

Inflation targets must constrain implemented rules; otherwise identical unrestricted rule classes make \(W(\bar\pi)\) independent of the label. For example restrict to a fixed family \(i_t=\max\{\underline i,\bar r+\bar\pi+\phi_\pi(\pi_t-\bar\pi)+\phi_xx_t\}\), with \((\phi_\pi,\phi_x)\) in a stated compact admissible set and rates and inflation in the same units. For parameter \(\theta\) in the population-compatible set \(\Theta(P)\), let \(W_\theta(\bar\pi)\) be the optimized welfare. The conditional optimum set is \(\bigcup_{\theta\in\Theta(P)}\arg\max_{\bar\pi}W_\theta(\bar\pi)\); a robust commitment instead maximizes \(\inf_{\theta\in\Theta(P)}W_\theta(\bar\pi)\), or minimizes specified welfare regret. These are different criteria. Compact targets do not guarantee attainment without nonempty rule feasibility and appropriate continuity or upper semicontinuity.
\subsection*{Variants and assumption changes}
The following are \emph{editorial variants}, not additional published open conjectures.
\begin{enumerate}
\item \emph{Heterogeneous inflation exposure.} Retain the same target/rule family but add household-specific nominal debt, wage indexing, and consumption baskets. Compare welfare-optimal targets under fixed social weights with the representative-household benchmark.
\item \emph{Minimax regret.} Choose a target minimizing \(\sup_{\theta\in\Theta(P)}[\max_{p}W_\theta(p)-W_\theta(\bar\pi)]\). This differs from worst-case welfare and allows comparisons when model utility levels have different arbitrary normalizations.
\end{enumerate}
\subsection*{References and source audit}
\begin{itemize}
\item Hess Chung; Callum Jones; Antoine Lepetit; Fernando M. Martin. \emph{Implications of Inflation Dynamics for Monetary Policy Strategies}. 2025. Locator: Abstract, printed p.1; Introduction, printed pp.2--3; strategy comparisons rather than choice of target level. Primary text checked. \url{https://www.federalreserve.gov/econres/feds/files/2025072pap.pdf}.
\end{itemize}
Strategy comparisons do not publish an unresolved optimal numerical target; the robust-target optimization is editorial.
\begin{enumerate}\raggedright
\item\hypertarget{definitions:source:30007528:1}{} Hess Chung; Callum Jones; Antoine Lepetit; Fernando M. Martin. 2025. \emph{Implications of Inflation Dynamics for Monetary Policy Strategies}. Abstract, printed p.1; Introduction, printed pp.2--3; strategy comparisons rather than choice of target level.
\url{https://www.federalreserve.gov/econres/feds/files/2025072pap.pdf}.
DOI: \url{https://doi.org/10.17016/FEDS.2025.072}.
\end{enumerate}

\subsection*{Common conventions: Source mathematical conventions}

These conventions apply to each entry unless explicitly replaced.

\textbf{Models and identification.} A primitive model \(m\) specifies all probability laws, feasible choices, information, equilibrium and measurement rules used by its target. The admissible class \(\mathcal M\) is fixed independently of outcomes. If \(P_O(m)\) is its observable probability law and \(\tau(m)\) the target, its identified set at observable law \(P\) is
\[
 \mathcal I_\tau(P)=\{\tau(m):m\in\mathcal M,\ P_O(m)=P\}.
\]
Point identification means that this set is a singleton. An empty set signals incompatibility of data and assumptions. Coordinate bounds need not describe the joint set. Bounds are sharp only if every retained target is attainable or arbitrarily approachable by compatible models. Finite bounds require explicit restrictions on unobserved quantities; unrestricted confounding, hidden wealth, extrapolation or arbitrary equilibrium selection can yield uninformative sets.

\textbf{Causal interventions.} A potential outcome \(Y_i(p)\) is the outcome under a complete policy regime \(p\), including timing, financing, information and market spillovers. The target population and units are fixed before treatment. With interference, use the complete assignment vector or a justified exposure mapping; individual no-interference notation alone cannot identify an economy-wide intervention. A difference of mean potential outcomes does not require the cross-world joint distribution. Conditioning on post-treatment migration, survival, enrollment or employment changes the estimand and needs separate justification.

\textbf{Statistical statements.} An estimator, confidence set, forecast or test specifies a sampling law, model class and loss/coverage criterion. Uniform \((1-\alpha)\)-coverage means \(\inf_{m\in\mathcal M}\Pr_m\{\tau(m)\in C_n\}\geq1-\alpha\), or an explicitly asymptotic version. The whole parameter space trivially has coverage, so informativeness requires a stated width, risk or power objective. A forecasting improvement is relative to a named benchmark, information set and evaluation population.

\textbf{Equilibrium and optimization.} An equilibrium comprises feasible plans, optimality, beliefs where needed, budget constraints and all market/resource clearing. The primitives or a stated benchmark define each correspondence \(\mathcal E(p;m)\). Nonemptiness is assumed or proved, never inferred just from compact policy sets. A selected-equilibrium target uses a declared selection rule; a robust target quantifies over all permitted equilibria. Use suprema or infima unless attainment is established. An \(\varepsilon\)-optimal policy is feasible and within \(\varepsilon\) of the finite optimum value. Report an empty feasible set separately.

\textbf{Welfare and accounting.} Normative utility, interpersonal normalization, social weights, numeraire, discounting and the use of public receipts are primitives. Behavioral utility need not coincide with the welfare criterion. Household welfare includes ownership income and rents once; adding profits again to compensating variation generally double-counts them. A monetary equivalent is defined by an explicit transfer and price convention, with monotonicity and range conditions. Average effects, mechanism contributions, transfers and welfare losses are different targets. Infinite sums require convergence and dynamic feasible plans require terminal or no-Ponzi conditions.

\textbf{Source versions.} A working-paper date, revision date and journal publication year are distinguished. A DOI for a journal version is not automatically the DOI of an earlier manuscript. Locators refer to the stated version and printed pages where specified. A metadata-only check does not verify a claim about a full-text conclusion. Every entry's source subsection narrows the provenance claim accordingly.
% END ECONOMICS STATEMENT
\end{document}
